% GENERATED FILE. Edit canonical lessons, then run npm run generate:books.
% canonical-source-hash: fnv1a64:4d68ca65d152df7c
% canonical-lessons: KA-C347-suddigosthi, KA-C347-prasnisu, KA-C347-pratikriyisu, KA-C347-ullekhisu, KA-C347-saramsa, KA-R347-review-341-347

\chapter{The Press Conference}
\label{ch:ka-the-press-conference}
\hlchaptermodality{347}

\begin{chapteropening}
\textbf{By the end of this chapter:} \emph{I can say a press conference, to question, to give a reaction, to mention and a summary.}

Complete the last lesson of chapter 347: I can say a press conference, to question, to give a reaction, to mention and a summary.
\end{chapteropening}

\section[\texorpdfstring{\kn{ಸುದ್ದಿಗೋಷ್ಠಿ} (suddigōṣṭhi)}{ಸುದ್ದಿಗೋಷ್ಠಿ (suddigōṣṭhi)}]{\kn{ಸುದ್ದಿಗೋಷ್ಠಿ} (suddigōṣṭhi) — a press conference}

\label{lesson:KA-C347-suddigosthi}



\begin{quote}
Before the new one: say the Kannada for a deluge, then the Kannada for an injured person.
\end{quote}

\subsection*{You'll want to know: \kn{ಸುದ್ದಿಗೋಷ್ಠಿ}}
\textbf{\kn{ಸುದ್ದಿಗೋಷ್ಠಿ}} — \emph{suddigōṣṭhi} — \textquotedblleft{}a press conference\textquotedblright{}.

\begin{quote}
\kn{ಪೊಲೀಸ್} \kn{ಅಧಿಕಾರಿ} \kn{ಸಂಜೆ} \kn{ಸುದ್ದಿಗೋಷ್ಠಿ} \kn{ಕರೆದರು}.
\end{quote}

\begin{quote}
\emph{polīs adhikāri sañje suddigōṣṭhi karedaru.}
\end{quote}

\begin{quote}
The police officer called a press conference in the evening.
\end{quote}

A press conference is \textquotedblleft{}called\textquotedblright{}: \kn{ಸುದ್ದಿಗೋಷ್ಠಿ} \kn{ಕರೆ}, or \textquotedblleft{}held\textquotedblright{}: \kn{ಸುದ್ದಿಗೋಷ್ಠಿ} \kn{ನಡೆಸು}.

\begin{cousinweb}
\textbf{\kn{ಸುದ್ದಿಗೋಷ್ಠಿ}} (\emph{suddigōṣṭhi}) is a Kannada compound of \kn{ಸುದ್ದಿ}, \textquotedblleft{}news\textquotedblright{}, which comes from Sanskrit \emph{śuddhi} (CDIAL 12523), and \kn{ಗೋಷ್ಠಿ}, borrowed from Sanskrit \emph{goṣṭhī}, \textquotedblleft{}an assembly, a meeting, a conversation\textquotedblright{} (Monier-Williams).
\end{cousinweb}

\begin{grammarlens}[title={What you've built}]
One more word for telling what happened.
\end{grammarlens}

\subsection*{Your turn}
\begin{itemize}
\raggedright
  \item \emph{Say it:} \emph{suddigōṣṭhi}
  \item \emph{Say it:} \emph{suddigōṣṭhi}, once more
  \item \emph{Say it:} \emph{polīs adhikāri sañje suddigōṣṭhi karedaru.} — The police officer called a press conference in the evening.
  \item \emph{Recall:} say the Kannada for a deluge, then the Kannada for an injured person, then say \emph{suddigōṣṭhi} again
\end{itemize}

\begin{tcolorbox}[breakable,colback=teal!4,colframe=teal!35!black,title={Before you move on}]
What does \kn{ಸುದ್ದಿಗೋಷ್ಠಿ} mean? (\textquotedblleft{}A press conference\textquotedblright{}.) Say it once more.
\end{tcolorbox}
\section[\texorpdfstring{\kn{ಪ್ರಶ್ನಿಸು} (praśnisu)}{ಪ್ರಶ್ನಿಸು (praśnisu)}]{\kn{ಪ್ರಶ್ನಿಸು} (praśnisu) — to question}

\label{lesson:KA-C347-prasnisu}



\begin{quote}
Before the new one: say the Kannada for an injured person, then the Kannada for a press conference.
\end{quote}

\subsection*{You'll want to know: \kn{ಪ್ರಶ್ನಿಸು}}
\textbf{\kn{ಪ್ರಶ್ನಿಸು}} — \emph{praśnisu} — \textquotedblleft{}to question\textquotedblright{}.

\begin{quote}
\kn{ವರದಿಗಾರರು} \kn{ಅಧಿಕಾರಿಯನ್ನು} \kn{ಅಪಘಾತದ} \kn{ಬಗ್ಗೆ} \kn{ಪ್ರಶ್ನಿಸಿದರು}.
\end{quote}

\begin{quote}
\emph{varadigāraru adhikāriyannu apaghātada bagge praśnisidaru.}
\end{quote}

\begin{quote}
The reporters questioned the officer about the accident.
\end{quote}

Past (I …): \kn{ಪ್ರಶ್ನಿಸಿದೆ}, \emph{praśniside}.

The person questioned takes -\kn{ಅನ್ನು}. It can also mean to challenge: \kn{ನಿರ್ಧಾರವನ್ನು} \kn{ಪ್ರಶ್ನಿಸು}.

\begin{cousinweb}
\textbf{\kn{ಪ್ರಶ್ನಿಸು}} (\emph{praśnisu}) is built on Sanskrit \emph{praśna}, \textquotedblleft{}a question\textquotedblright{} (CDIAL 8818), with the Kannada verb ending -\emph{isu}.
\end{cousinweb}

\begin{grammarlens}[title={What you've built}]
One more word for telling what happened.
\end{grammarlens}

\subsection*{Your turn}
\begin{itemize}
\raggedright
  \item \emph{Say it:} \emph{praśnisu}
  \item \emph{Say it:} \emph{praśnisu}, once more
  \item \emph{Say it:} \emph{varadigāraru adhikāriyannu apaghātada bagge praśnisidaru.} — The reporters questioned the officer about the accident.
  \item \emph{Recall:} say the Kannada for an injured person, then the Kannada for a press conference, then say \emph{praśnisu} again
\end{itemize}

\begin{tcolorbox}[breakable,colback=teal!4,colframe=teal!35!black,title={Before you move on}]
What does \kn{ಪ್ರಶ್ನಿಸು} mean? (\textquotedblleft{}To question\textquotedblright{}.) Say it once more.
\end{tcolorbox}
\section[\texorpdfstring{\kn{ಪ್ರತಿಕ್ರಿಯಿಸು} (pratikriyisu)}{ಪ್ರತಿಕ್ರಿಯಿಸು (pratikriyisu)}]{\kn{ಪ್ರತಿಕ್ರಿಯಿಸು} (pratikriyisu) — to give a reaction; to respond}

\label{lesson:KA-C347-pratikriyisu}



\begin{quote}
Before the new one: say the Kannada for a press conference, then the Kannada for to question.
\end{quote}

\subsection*{You'll want to know: \kn{ಪ್ರತಿಕ್ರಿಯಿಸು}}
\textbf{\kn{ಪ್ರತಿಕ್ರಿಯಿಸು}} — \emph{pratikriyisu} — \textquotedblleft{}to give a reaction; to respond\textquotedblright{}.

\begin{quote}
\kn{ಆ} \kn{ಪ್ರಶ್ನೆಗೆ} \kn{ಅಧಿಕಾರಿ} \kn{ಪ್ರತಿಕ್ರಿಯಿಸಲಿಲ್ಲ}.
\end{quote}

\begin{quote}
\emph{ā praśnege adhikāri pratikriyisalilla.}
\end{quote}

\begin{quote}
The officer did not respond to that question.
\end{quote}

Past (I …): \kn{ಪ್ರತಿಕ್ರಿಯಿಸಿದೆ}, \emph{pratikriyiside}.

What one responds to takes -\kn{ಗೆ}. The noun is \kn{ಪ್ರತಿಕ್ರಿಯೆ}, a response.

\begin{cousinweb}
\textbf{\kn{ಪ್ರತಿಕ್ರಿಯಿಸು}} (\emph{pratikriyisu}) is built on Sanskrit \emph{pratikriyā}, \textquotedblleft{}requital, counteraction\textquotedblright{} (Monier-Williams), from \emph{prati} \textquotedblleft{}back, in return\textquotedblright{} and \emph{kriyā} \textquotedblleft{}an act\textquotedblright{}, with the Kannada verb ending -\emph{isu}.
\end{cousinweb}

\begin{grammarlens}[title={What you've built}]
One more word for telling what happened.
\end{grammarlens}

\subsection*{Your turn}
\begin{itemize}
\raggedright
  \item \emph{Say it:} \emph{pratikriyisu}
  \item \emph{Say it:} \emph{pratikriyisu}, once more
  \item \emph{Say it:} \emph{ā praśnege adhikāri pratikriyisalilla.} — The officer did not respond to that question.
  \item \emph{Recall:} say the Kannada for a press conference, then the Kannada for to question, then say \emph{pratikriyisu} again
\end{itemize}

\begin{tcolorbox}[breakable,colback=teal!4,colframe=teal!35!black,title={Before you move on}]
What does \kn{ಪ್ರತಿಕ್ರಿಯಿಸು} mean? (\textquotedblleft{}To give a reaction; to respond\textquotedblright{}.) Say it once more.
\end{tcolorbox}
\section[\texorpdfstring{\kn{ಉಲ್ಲೇಖಿಸು} (ullēkhisu)}{ಉಲ್ಲೇಖಿಸು (ullēkhisu)}]{\kn{ಉಲ್ಲೇಖಿಸು} (ullēkhisu) — to mention; to refer to}

\label{lesson:KA-C347-ullekhisu}



\begin{quote}
Before the new one: say the Kannada for to question, then the Kannada for to give a reaction.
\end{quote}

\subsection*{You'll want to know: \kn{ಉಲ್ಲೇಖಿಸು}}
\textbf{\kn{ಉಲ್ಲೇಖಿಸು}} — \emph{ullēkhisu} — \textquotedblleft{}to mention; to refer to\textquotedblright{}.

\begin{quote}
\kn{ಅವರು} \kn{ಮಾತಿನಲ್ಲಿ} \kn{ನಿನ್ನೆಯ} \kn{ಘಟನೆಯನ್ನು} \kn{ಉಲ್ಲೇಖಿಸಿದರು}.
\end{quote}

\begin{quote}
\emph{avaru mātinalli ninneya ghaṭaneyannu ullēkhisidaru.}
\end{quote}

\begin{quote}
In the talk he mentioned yesterday's incident.
\end{quote}

Past (I …): \kn{ಉಲ್ಲೇಖಿಸಿದೆ}, \emph{ullēkhiside}.

What is mentioned takes -\kn{ಅನ್ನು}. A formal word; in talk people say \kn{ಹೇಳು}.

\begin{cousinweb}
\textbf{\kn{ಉಲ್ಲೇಖಿಸು}} (\emph{ullēkhisu}) is built on Sanskrit \emph{ullekha}, \textquotedblleft{}causing to come forth clearly\textquotedblright{} (Monier-Williams), with the Kannada verb ending -\emph{isu}.
\end{cousinweb}

\begin{grammarlens}[title={What you've built}]
One more word for telling what happened.
\end{grammarlens}

\subsection*{Your turn}
\begin{itemize}
\raggedright
  \item \emph{Say it:} \emph{ullēkhisu}
  \item \emph{Say it:} \emph{ullēkhisu}, once more
  \item \emph{Say it:} \emph{avaru mātinalli ninneya ghaṭaneyannu ullēkhisidaru.} — In the talk he mentioned yesterday's incident.
  \item \emph{Recall:} say the Kannada for to question, then the Kannada for to give a reaction, then say \emph{ullēkhisu} again
\end{itemize}

\begin{tcolorbox}[breakable,colback=teal!4,colframe=teal!35!black,title={Before you move on}]
What does \kn{ಉಲ್ಲೇಖಿಸು} mean? (\textquotedblleft{}To mention; to refer to\textquotedblright{}.) Say it once more.
\end{tcolorbox}
\section[\texorpdfstring{\kn{ಸಾರಾಂಶ} (sārāṃśa)}{ಸಾರಾಂಶ (sārāṃśa)}]{\kn{ಸಾರಾಂಶ} (sārāṃśa) — a summary; the gist}

\label{lesson:KA-C347-saramsa}



\begin{quote}
Before the new one: say the Kannada for to give a reaction, then the Kannada for to mention.
\end{quote}

\subsection*{You'll want to know: \kn{ಸಾರಾಂಶ}}
\textbf{\kn{ಸಾರಾಂಶ}} — \emph{sārāṃśa} — \textquotedblleft{}a summary; the gist\textquotedblright{}.

\begin{quote}
\kn{ಸುದ್ದಿಗೋಷ್ಠಿಯ} \kn{ಸಾರಾಂಶ} \kn{ರಾತ್ರಿಯ} \kn{ವಾರ್ತೆಯಲ್ಲಿ} \kn{ಬಂತು}.
\end{quote}

\begin{quote}
\emph{suddigōṣṭhiya sārāṃśa rātriya vārteyalli bantu.}
\end{quote}

\begin{quote}
A summary of the press conference came on the night news.
\end{quote}

\kn{ಸಾರ} is the essence: the essential part of something said or written.

\begin{cousinweb}
\textbf{\kn{ಸಾರಾಂಶ}} (\emph{sārāṃśa}) is borrowed from Sanskrit, a compound of \emph{sāra}, \textquotedblleft{}the core, the essence\textquotedblright{}, and \emph{aṃśa}, \textquotedblleft{}a part, a share\textquotedblright{} (Monier-Williams).
\end{cousinweb}

\begin{grammarlens}[title={What you've built}]
One more word for telling what happened.
\end{grammarlens}

\subsection*{Your turn}
\begin{itemize}
\raggedright
  \item \emph{Say it:} \emph{sārāṃśa}
  \item \emph{Say it:} \emph{sārāṃśa}, once more
  \item \emph{Say it:} \emph{suddigōṣṭhiya sārāṃśa rātriya vārteyalli bantu.} — A summary of the press conference came on the night news.
  \item \emph{Recall:} say the Kannada for to give a reaction, then the Kannada for to mention, then say \emph{sārāṃśa} again
\end{itemize}

\begin{tcolorbox}[breakable,colback=teal!4,colframe=teal!35!black,title={Before you move on}]
What does \kn{ಸಾರಾಂಶ} mean? (\textquotedblleft{}A summary; the gist\textquotedblright{}.) Say it once more.
\end{tcolorbox}
\section[(dialogue)]{Looking back: from Our Family Line to The Press Conference}

\label{lesson:KA-R347-review-341-347}



\begin{quote}
Say the Kannada for a lineage and an ancestor.
\end{quote}

\subsection*{Your turn}
Say these aloud:

\begin{itemize}
\raggedright
  \item a lineage, an ancestor, a family generation, a family line, a monarch
  \item to besiege, a moat, to surround, to surrender, a cannon
  \item to narrate, to describe vividly, a situation, an episode, a background
  \item a weapon, a gun, a battlefield, a chariot, to do battle
  \item to hide oneself, to go into hiding, to chase, a spy, to sneak in
  \item an explosion, to reel, to tremble, a deluge, an injured person
  \item a press conference, to question, to give a reaction, to mention, a summary
\end{itemize}

\begin{tcolorbox}[breakable,colback=teal!4,colframe=teal!35!black,title={Before you move on}]
A lineage? (\textbf{\kn{ವಂಶ}}, \emph{vaṃśa}.) An ancestor? (\textbf{\kn{ಪೂರ್ವಜ}}, \emph{pūrvaja}.) A family generation? (\textbf{\kn{ತಲೆಮಾರು}}, \emph{talemāru}.) A family line? (\textbf{\kn{ಮನೆತನ}}, \emph{manetana}.) A monarch? (\textbf{\kn{ಅರಸ}}, \emph{arasa}.) To besiege? (\textbf{\kn{ಮುತ್ತಿಗೆ} \kn{ಹಾಕು}}, \emph{muttige hāku}.) A moat? (\textbf{\kn{ಕಂದಕ}}, \emph{kandaka}.) To surround? (\textbf{\kn{ಸುತ್ತುವರಿ}}, \emph{suttuvari}.) To surrender? (\textbf{\kn{ಶರಣಾಗು}}, \emph{śaraṇāgu}.) A cannon? (\textbf{\kn{ಫಿರಂಗಿ}}, \emph{phiraṅgi}.) To narrate? (\textbf{\kn{ನಿರೂಪಿಸು}}, \emph{nirūpisu}.) To describe vividly? (\textbf{\kn{ವರ್ಣಿಸು}}, \emph{varṇisu}.) A situation? (\textbf{\kn{ಸನ್ನಿವೇಶ}}, \emph{sannivēśa}.) An episode? (\textbf{\kn{ಪ್ರಸಂಗ}}, \emph{prasaṅga}.) A background? (\textbf{\kn{ಹಿನ್ನೆಲೆ}}, \emph{hinnele}.) A weapon? (\textbf{\kn{ಆಯುಧ}}, \emph{āyudha}.) A gun? (\textbf{\kn{ಬಂದೂಕು}}, \emph{bandūku}.) A battlefield? (\textbf{\kn{ರಣರಂಗ}}, \emph{raṇaraṅga}.) A chariot? (\textbf{\kn{ರಥ}}, \emph{ratha}.) To do battle? (\textbf{\kn{ಕಾದಾಡು}}, \emph{kādāḍu}.) To hide oneself? (\textbf{\kn{ಅಡಗು}}, \emph{aḍagu}.) To go into hiding? (\textbf{\kn{ತಲೆಮರೆಸಿಕೊ}}, \emph{talemaresiko}.) To chase? (\textbf{\kn{ಬೆನ್ನಟ್ಟು}}, \emph{bennaṭṭu}.) A spy? (\textbf{\kn{ಗೂಢಚಾರ}}, \emph{gūḍhacāra}.) To sneak in? (\textbf{\kn{ನುಸುಳು}}, \emph{nusuḷu}.) An explosion? (\textbf{\kn{ಸ್ಫೋಟ}}, \emph{sphōṭa}.) To reel? (\textbf{\kn{ತತ್ತರಿಸು}}, \emph{tattarisu}.) To tremble? (\textbf{\kn{ಕಂಪಿಸು}}, \emph{kampisu}.) A deluge? (\textbf{\kn{ಪ್ರಳಯ}}, \emph{praḷaya}.) An injured person? (\textbf{\kn{ಗಾಯಾಳು}}, \emph{gāyāḷu}.) A press conference? (\textbf{\kn{ಸುದ್ದಿಗೋಷ್ಠಿ}}, \emph{suddigōṣṭhi}.) To question? (\textbf{\kn{ಪ್ರಶ್ನಿಸು}}, \emph{praśnisu}.) To give a reaction? (\textbf{\kn{ಪ್ರತಿಕ್ರಿಯಿಸು}}, \emph{pratikriyisu}.) To mention? (\textbf{\kn{ಉಲ್ಲೇಖಿಸು}}, \emph{ullēkhisu}.) A summary? (\textbf{\kn{ಸಾರಾಂಶ}}, \emph{sārāṃśa}.)
\end{tcolorbox}
